\section{Partial Fine-tuning}
% background + motivation
This type of tuning methods (which is also called \textit{specification-based tuning}, \citealp{DBLP:journals/corr/abs-2203-06904}) fine-tune a few inherent parameters while leaving the majority of parameters unchanged in model adaptation. This approach does not seek to change the internal structure of a model but to optimize a small number of internal parameters to solve particular tasks. Generally, such specifications could be implemented based on heuristics or training supervision. 

\paragraph{Heuristic Specification.}
Specification-based methods do not introduce any new parameters in the model, but directly specify part of the parameters to be optimized. The idea is simple but surprisingly effective, \citet{DBLP:journals/corr/abs-1911-03090} \textbf{only fine-tune one-fourth of the final layers} of BERT and RoBERTa and could produce $90\%$ of the performance of full parameter fine-tuning. 

BitFit~\citep{DBLP:conf/acl/ZakenGR22} empirically proves that by \textbf{only optimizing the bias terms} inside the model and freezing other parameters, the model could still reproduce over $95\%$ performance on several benchmarks. 
The biased term exists in both attention mechanism and MLP feedforward layer. BitFit can fine-tune only two bias components (the “query” and “middle-of-MLP” bias terms), amounting to half of the bias parameters in the model, and only $0.04\%$ of all model parameters. Corresponding bias terms are colored red. Other bias terms are colored in blue.
\begin{align}
    \texttt{Query in Attention}:& \ \ \textbf{Q}(\textbf{x}) = \textbf{W}_{q}\textbf{x} + {\color{red}{\textbf{b}_{q}}} \\
    & \ \ \textbf{K}(\textbf{x}) = \textbf{W}_{k}\textbf{x} + {\color{blue}{\textbf{b}_{k}}} \\
    & \ \ \textbf{V}(\textbf{x}) = \textbf{W}_{v}\textbf{x} + {\color{blue}{\textbf{b}_{v}}} \\
    \texttt{MLP in Feedfoward}:& \ \ \textbf{h}_{2} = \text{Dropout}(\textbf{W}_{1} \cdot \textbf{h}_{1} + {\color{blue}{\textbf{b}_{1}}}) \\
    & \ \ \textbf{h}_{3} = \textbf{g}_{LN_{1}} \odot \frac{(\textbf{h}_{2}+\textbf{x})-\mu}{\sigma} + {\color{blue}{\textbf{b}_{LN_{1}}}} \\
    & \ \ \textbf{h}_{4} = \text{GELU}(\textbf{W}_{2} \cdot \textbf{h}_{3} + {\color{red}{\textbf{b}_{2}}}) \\
    & \ \ \textbf{h}_{5} = \text{Dropout}(\textbf{W}_{3} \cdot \textbf{h}_{4} + {\color{blue}{\textbf{b}_{3}}}) \\
    & \ \ \text{out} = \textbf{g}_{LN_{2}} \odot \frac{(\textbf{h}_{5}+\textbf{h}_{3})-\mu}{\sigma} + {\color{blue}{\textbf{b}_{LN_{2}}}} 
\end{align}
Empirical results in BitFit also show that even if we use a small random set of parameters for tuning (which obviously will degrade the performance), the model could still yield passable results on the GLUE benchmark. 
Unfortunately, the work only applies this trick to small-scale models, and there is no guarantee that randomly choosing some parameters to be tuned would remain competitive for larger models. Another valuable observation is that different bias terms may have different functionalities during model adaptation.

\paragraph{Learn the Specification.}
Rather than manually or heuristically specify which parameters to be updated, one alternative is to ``learn'' such specifications. \textit{Diff pruning}~\citep{DBLP:conf/acl/GuoRK20} reparameterizes the fine-tuned model parameters as the summation of the pre-trained parameters and the difference vector as $\theta_{FT} = \theta_{LM} + \delta_{Diff} $. Hence, the key issue is to encourage the difference vector $\delta_{Diff}$ to be as sparse as possible. This work regularizes the vector by a differentiable approximation to the $L_0$-norm penalty as $||\delta_{Diff}||_{0}$ to achieve the goal of sparsity. Practically, because new parameters to be optimized are introduced in the learning phase, \textit{Diff pruning} takes up more GPU memory than full parameter fine-tuning, which may establish barriers in the application on large language models. 

%\mrinmaya{remove this; the details are not clear from this} The masking method~\citep{DBLP:conf/emnlp/ZhaoLMJS20} learns selective masks for language models to only update the critical weights for particular tasks. To learn such a set of masks, a binary matrix associated with the model weights is introduced, where each value is generated by a thresholding function. During back-propagation, the matrix is updated by a noisy estimator.

